\documentclass[./main.tex]{subfiles}

\begin{document}

\begin{center}
    {\Large\bfseries What Uses All the Energy?}\\[4pt]
    A Survey of Manufacturing Facility Systems
\end{center}

\section*{Model 1: The Energy Landscape}

Different types of manufacturing facilities use energy in very different ways.
The table below shows energy end-use breakdowns for four facilities, each in a
different industrial sector. Note that these values are specific to these four
facilities. Energy allocation varies widely between facilities, even within the
same sector, depending on specific processes, equipment age, climate zone, and
operating schedule.

\begin{longtable}[]{@{}
  >{\raggedright\arraybackslash}p{(\columnwidth - 8\tabcolsep) * \real{0.2350}}
  >{\raggedright\arraybackslash}p{(\columnwidth - 8\tabcolsep) * \real{0.1912}}
  >{\raggedright\arraybackslash}p{(\columnwidth - 8\tabcolsep) * \real{0.1912}}
  >{\raggedright\arraybackslash}p{(\columnwidth - 8\tabcolsep) * \real{0.1912}}
  >{\raggedright\arraybackslash}p{(\columnwidth - 8\tabcolsep) * \real{0.1912}}@{}}
\toprule\noalign{}
Energy end use & Facility A: Food Processing & Facility B: Metal Fabrication
  & Facility C: Chemical Mfg. & Facility D: Auto Assembly \\
\midrule\noalign{}
\endhead
\bottomrule\noalign{}
\endlastfoot
Motors / drives   & 25\% & 45\% & 50\% & 40\% \\
Process heating   & 35\% & 15\% & 20\% & 10\% \\
HVAC              & 13\% &  7\% &  5\% & 20\% \\
Compressed air    &  5\% & 18\% & 10\% & 15\% \\
Refrigeration     & 12\% &  2\% &  3\% &  2\% \\
Lighting          &  5\% &  6\% &  4\% &  6\% \\
Other             &  5\% &  7\% &  8\% &  7\% \\
\end{longtable}

\emph{Table 1. Energy end-use shares (\%) for four manufacturing facilities in
different industrial sectors. These are facility-specific values; other
facilities in the same sector may differ significantly.}

\begin{enumerate}
  \needspace{4\baselineskip}\item Which energy end-use category claims the largest share in three of the
        four facilities shown?
        \vspace{1.25in}

  \needspace{4\baselineskip}\item Which facility is the exception to the pattern you identified in
        Question 1? What end-use category dominates there, and why might that
        be the case given what that facility produces?
        \vspace{1.25in}

  \needspace{4\baselineskip}\item Refrigeration accounts for 12\% of energy use at Facility A (food
        processing) but only 2--3\% at the other three. In one sentence,
        explain why this might be.
        \vspace{1.25in}

  \needspace{4\baselineskip}\item Rank the four facilities from highest to lowest HVAC energy share.
        What characteristic of Facility A (food processing) and Facility D
        (automotive assembly) might explain why HVAC is a much larger energy
        consumer there than the others?
        \vspace{1.25in}

  \needspace{4\baselineskip}\item Suppose you read an industry report claiming that ``motors typically
        account for about 40\% of energy use in manufacturing.'' Based on
        Table 1, is this a reasonable generalization? What important caveat
        would you add?
        \vspace{1.25in}

  \needspace{4\baselineskip}\item A facility manager at Facility C (chemical manufacturing) tells you,
        ``We don't need to worry about compressed air; it's only 10\% of our
        bill.'' Without doing any calculations, give two reasons why this
        statement might be misleading.
        \vspace{1.25in}
\end{enumerate}

\clearpage
\section*{Model 2: A Field Guide to Facility Systems}

The table below introduces the major energy-consuming system types you will
encounter during an industrial energy assessment. For each system, a single key
performance metric is listed. This is the number that assessors most commonly
use as a starting point when evaluating how well the system is performing.
These metrics are not the only way to evaluate each system, but they give you a
consistent vocabulary for comparing equipment and identifying where to look
more closely.

\begin{longtable}[]{@{}
  >{\raggedright\arraybackslash}p{(\columnwidth - 4\tabcolsep) * \real{0.1923}}
  >{\raggedright\arraybackslash}p{(\columnwidth - 4\tabcolsep) * \real{0.4274}}
  >{\raggedright\arraybackslash}p{(\columnwidth - 4\tabcolsep) * \real{0.3803}}@{}}
\toprule\noalign{}
System & What it does & Key metric \\
\midrule\noalign{}
\endhead
\bottomrule\noalign{}
\endlastfoot
Lighting          & Provides illumination for tasks, safety, and general
                    visibility throughout the facility
                  & Lumens per watt (lm/W) \\
Motors / drives   & Convert electrical energy into mechanical energy to power
                    fans, pumps, conveyors, and production equipment
                  & \% nominal efficiency (full-load) \\
Compressed air    & Generate and distribute pressurized air for pneumatic
                    tools, actuators, blow-off, and process applications
                  & kW/100 cfm (specific power) \\
HVAC              & Maintain thermal comfort and air quality in occupied and
                    conditioned spaces
                  & kW/ton (cooling), \% AFUE (heating) \\
Process heating   & Apply thermal energy directly to materials for drying,
                    curing, melting, heat treating, or other transformation
                  & \% combustion efficiency \\
Refrigeration     & Remove heat from products or spaces to maintain
                    temperatures below ambient for storage or process cooling
                  & kW/ton \\
Steam systems     & Generate and distribute steam for process heating,
                    mechanical drives, or building heat
                  & \% combustion efficiency \\
\end{longtable}

\emph{Table 2. Major energy system types encountered during industrial energy
assessments, with primary function and key performance metric.}

\begin{enumerate}
  \setcounter{enumi}{6}
  \needspace{4\baselineskip}\item Look at the key metrics column. Most of these metrics are ratios. What
        do the numerator and denominator of each ratio usually represent?
        \vspace{1.25in}

  \needspace{4\baselineskip}\item Sort the eight systems into two groups: those whose primary energy
        input is electricity, and those whose primary energy input is fuel
        (natural gas, oil, etc.). Identify at least one system that could
        reasonably fall into either category.
        \vspace{1.25in}

  \needspace{4\baselineskip}\item HVAC and refrigeration both involve removing heat from a space. Based
        on Table 2, what distinguishes them from each other in terms of
        purpose?
        \vspace{1.25in}

\end{enumerate}

\clearpage
\section*{Model 3: Component Efficiency vs.\ System Efficiency}

When you evaluate a piece of equipment, the nameplate or catalog efficiency
tells you how well that individual component converts its energy input to
useful output under rated conditions. But in practice, energy is lost at every
stage between the meter and the point of use. The diagram below illustrates
this concept for a compressed air system, tracing electrical energy from the
utility meter to the point where air does useful work.

\begin{figure}[ht]
  \centering
  \fbox{\parbox{0.85\textwidth}{\centering\vspace{2.5cm}
    \textit{[Figure 1: Compressed air system energy flow diagram --- insert image here]}
  \vspace{2.5cm}}}
  \caption{Compressed air system energy flow from utility meter to point of
    use, illustrating cumulative losses at each stage.}
\end{figure}

Starting from 100 kW of electrical input at the utility meter, motor losses
consume roughly 5\%, leaving about 95 kW of shaft power. The compression
process converts most of that mechanical energy into heat; only about 19 kW
remains as useful air energy at the compressor discharge. Distribution losses
through piping (pressure drop, friction) reduce this to roughly 16 kW at the
point of use. Finally, leaks in a typical system waste about 25\% of the
generated air, leaving approximately 12 kW of equivalent useful work. The
compressor motor itself is about 95\% efficient. The system as a whole is about
12\% efficient.

The same pattern, with different magnitudes, appears in other systems. A boiler
might be 82\% efficient at converting fuel to steam, but the steam system as a
whole (including distribution losses, trap failures, condensate return, and
insulation deficiencies) might deliver only 55--60\% of the original fuel
energy to the process.

\begin{enumerate}
  \setcounter{enumi}{11}
  \needspace{4\baselineskip}\item The compressor motor in Figure 1 is approximately 95\% efficient. If
        you were evaluating this system by component efficiency alone, would it
        appear to need attention? Why or why not?
        \vspace{1.25in}

  \needspace{4\baselineskip}\item Where in the system does the single largest energy loss occur? Is this
        loss something you could see or hear during a walk-through?
        \vspace{1.25in}

  \needspace{4\baselineskip}\item Leaks account for roughly 25\% of the generated air in this example.
        If you could eliminate all leaks, how much of the original 100 kW
        input would reach useful work? Show your reasoning.
        \vspace{1.25in}

  \needspace{4\baselineskip}\item In your own words, explain the difference between component efficiency
        and system efficiency. Why does this distinction matter for energy
        assessors?
        \vspace{1.25in}

  \needspace{4\baselineskip}\item A facility manager says, ``We just bought a new premium-efficiency
        compressor motor, so our compressed air system is now efficient.''
        Based on Figure 1, explain why this claim is incomplete.
        \vspace{1.25in}

  \needspace{4\baselineskip}\item Name one other system from Table 2 where you would expect a large gap
        between component efficiency and system efficiency. Briefly explain
        what losses occur between the prime mover and the useful output.
        \vspace{1.25in}
\end{enumerate}

\clearpage
\section*{Model 4: Spotting Opportunities on a Walk-Through}

Below are ten observations recorded by an assessment team during a walk-through
of a medium-sized metals fabrication plant. The facility operates one shift
(6 AM to 2:30 PM, Monday through Friday) and has approximately 120,000 square
feet of conditioned production space plus 30,000 square feet of unconditioned
warehouse.

Read each observation and use what you have learned from Models 1 through 3 to
answer the questions.

\begin{longtable}[]{@{}
  >{\raggedright\arraybackslash}p{(\columnwidth - 2\tabcolsep) * \real{0.0641}}
  >{\raggedright\arraybackslash}p{(\columnwidth - 2\tabcolsep) * \real{0.9359}}@{}}
\toprule\noalign{}
\textbf{\#} & \textbf{Observation} \\
\midrule\noalign{}
\endhead
\bottomrule\noalign{}
\endlastfoot
\textbf{A} & High-bay HID fixtures (400W metal halide) are on throughout the
  production floor, including areas where no work is being performed during the
  walk-through. \\
\textbf{B} & The main air compressor (75 HP rotary screw) runs continuously,
  including during the lunch break when most pneumatic tools are idle. \\
\textbf{C} & Several overhead doors between the conditioned production area and
  the unconditioned warehouse are propped open. \\
\textbf{D} & A 50 HP exhaust fan on the roof runs at full speed 24/7, though it
  is only needed during production hours for weld fume extraction. \\
\textbf{E} & Two audible air leaks are noticed near quick-disconnect fittings
  on the main compressed air header. \\
\textbf{F} & The natural gas furnace used for heat treating has visible flue
  gases with a slight yellow tinge, and the facility has no record of a recent
  combustion efficiency test. \\
\textbf{G} & The roof of the warehouse section has minimal insulation (estimated
  R-5), and the HVAC system for the adjacent production area runs year-round. \\
\textbf{H} & Several 25 HP motors driving conveyor systems are original
  equipment from the 1990s and are rewound, not replaced. \\
\textbf{I} & The chilled water system (80-ton reciprocating chiller) operates at
  full capacity despite the outdoor temperature being 55\textdegree F during
  the assessment. \\
\textbf{J} & Fluorescent T12 fixtures with magnetic ballasts are still in use
  in the office and breakroom areas. \\
\end{longtable}

\emph{Table 3. Walk-through observations from a metals fabrication facility
operating one shift, Monday through Friday, 6:00 AM to 2:30 PM.}

\begin{enumerate}
  \setcounter{enumi}{17}
  \needspace{4\baselineskip}\item For each observation (A through J), identify which system category from
        Table 2 it belongs to. Some observations may relate to more than one
        system.
        \vspace{1.25in}

  \needspace{4\baselineskip}\item Observations B and E both involve the compressed air system. Using what
        you learned in Model 3, explain how these two observations are related
        and why addressing them together would produce a larger savings than
        addressing either one alone.
        \vspace{1.25in}

  \needspace{4\baselineskip}\item Observation D describes a motor running 24/7 when it is only needed
        during one shift. Without doing a precise calculation, estimate roughly
        what fraction of this fan's annual energy consumption is wasted.
        (\emph{Hint: think about how many hours are in a year versus how many
        hours the fan is actually needed.})
        \vspace{1.25in}

  \needspace{4\baselineskip}\item Look at your answers to Question 18 and refer back to Table 1 (energy
        end-use shares for metals fabrication). Which three observations would
        you prioritize for further investigation, and why? Your answer should
        consider both the energy share of the relevant system category and the
        apparent severity of the issue observed.
        \vspace{1.25in}

  \needspace{4\baselineskip}\item Observation H mentions motors that have been rewound rather than
        replaced. Using Table 2's key metric for motors, explain why a rewound
        motor might perform differently than a new replacement. What additional
        information would you need to quantify the potential savings?
        \vspace{1.25in}

  \needspace{4\baselineskip}\item For one observation of your choice, describe what additional
        measurements or data you would collect before writing an assessment
        recommendation. Be specific: what instrument or data source, and what
        value are you trying to determine?
        \vspace{1.25in}
\end{enumerate}

\end{document}
